% ============================================================
% 01_introduction.tex — Introduction (Macedonian)
% ============================================================

\section{Контекст и мотивација}

Фантази Премиер Лига (Fantasy Premier League, FPL) е официјалната онлајн фантази игра поврзана со Англиската Премиер Лига, со повеќе од десет милиони активни менаџери ширум светот во сезоната 2024--25. Секој учесник составува виртуелен тим од петнаесет играчи во рамките на буџет од 100 милиони фунти, а поените се акумулираат врз основа на реалните перформанси на избраните играчи во секоја натпреварувачка недела. Успехот во FPL бара непрекинато следење на статистики, повреди, форма, распоред на натпреварувања и пазарни движења, па поради тоа управувањето со тимот во текот на цела сезона претставува сложен и напорен когнитивен процес.

Основниот предизвик во FPL е предвидливоста на перформансите на играчите. За разлика од традиционалните спортски натпревари каде победата или поразот е единствен исход, FPL резултатите се детерминирани од континуирани и високо варијабилни настани: голови, асистенции, чисти листови, пенали, жолти картони и минути на игра. Играчот кој во минатата недела постигна осум поени може следата недела да заработи нула поени поради повреда или ротација. Оваа случајна природа го прави проблемот особено атрактивен за методите на машинско учење, кои можат да идентификуваат статистички образци во историски податоци и да ги вградат во предикциски модели.

Покрај предвидувањето на индивидуалните перформанси, изборот на тим подлежи на строги правила: буџетско ограничување, позициска структура (2 голмани, 5 одбранбени играчи, 5 средни играчи и 3 напаѓачи), лимит од максимум три играчи по клуб, и систем на трансфери каде дополнителните промени носат казни. Ваквата структура на ограничувања го прави проблемот одличен кандидат за техниките на целобројно линеарно програмирање (ILP), кои обезбедуваат глобален оптимум во рамките на зададениот простор на решенија.

Во последните години, напредокот на јазичните модели (Large Language Models, LLMs) отвори нова димензија во автоматизираното одлучување. LLM-ите можат да толкуваат неструктурирани текстуални информации — извештаи, прес-конференции, новински написи — и да ги конвертираат во структурирани сигнали кои можат да ги надополнат квантитативните предикции. Претренирањето на LLM-ите на огромен обем текстови им дава способност за логичко расудување кое е тешко да се постигне со класичните ML методи.

\section{Изјава за проблемот}

Целта на оваа дипломска работа е да се проектира и имплементира целосно автоматизиран систем кој управува со FPL тим низ цела сезона без рачна интервенција. Формално, системот треба да:

\begin{itemize}
  \item Предвидува очекувани поени за секој играч пред секоја натпреварувачка недела.
  \item Избира оптимален состав на 15 играчи во рамките на сите FPL ограничувања.
  \item Одлучува за трансфери, капитенски избор и активирање на специјалните чипови (Triple Captain, Bench Boost, Wildcard, Free Hit).
  \item Интегрира информации во реално време за повреди и ротации пред секој рок.
  \item Обезбедува интерпретабилни нарације за секоја одлука.
\end{itemize}

Системот е евалуиран на сезоната 2025--26, почнувајќи од натпреварувачка недела 1 (GW1) со нулта информација за тековната сезона (\textit{blind test}), при користење само на историски податоци до сезоната 2024--25.

\section{Придонеси на тезата}

Оваа дипломска работа дава неколку оригинални придонеси:

\begin{enumerate}
  \item \textbf{Тројно-слојна хибридна архитектура}: Прв систем кој ги комбинира ML предикциите (LightGBM), ILP оптимизацијата (PuLP) и LLM интелигенцијата (Claude API, Gemini API) во единствена пипелајн за управување со FPL тим.

  \item \textbf{Pre-deadline Intel Pipeline}: Систем од седум модули за собирање на информации во реално време (достапност, ризик од ротација, прес-конференции, пазарни сигнали) кои автоматски ги коригираат предикциите пред секој рок.

  \item \textbf{Сезонски симулатор со хиперпараметарска оптимизација}: Имплементација на joint random search (250 обиди) и Bayesian оптимизација со Optuna TPE (100 обиди) за истовремена оптимизација на ML хиперпараметрите и симулаторскиот конфигурациски простор.

  \item \textbf{Идентификација и поправка на критични грешки}: Во текот на развојот беа идентификувани и поправени пет критични програматски грешки, чиe откривање директно придонело за подобрување на резултатите и за разбирање на сложеноста на системот.
\end{enumerate}

\section{Структура на тезата}

Тезата е организирана во шест поглавја:

\textbf{Поглавје 2 — Преглед на литература} дава сеопфатен преглед на релевантните трудови во областа на gradient boosting, спортска аналитика, оптимизација, Bayesian оптимизација на хиперпараметри, фантази спорт и употребата на LLM-и во поддршката на одлучувањето. Идентификувана е критична празнина во литературата: ниеден постоечки труд не комбинира ML, ILP и LLM во единствена пипелајн за управување со FPL сезона.

\textbf{Поглавје 3 — Методологија} дава детален опис на целата системска архитектура, вклучувајќи ги собирањето и инженерингот на карактеристиките, тренирањето на LightGBM моделите, формулацијата на ILP проблемот, intel pipeline модулите, механизмот за оптимизација на хиперпараметрите и нарацискиот слој преку Claude API.

\textbf{Поглавје 4 — Резултати и евалуација} ги претставува квантитативните резултати од сезоната 2025--26 (GW1--28), историјата на подобрување на системот, влијанието на intel pipeline-от, анализата на активирањето на чиповите, и детален опис на грешките кои беа пронајдени и поправени.

\textbf{Поглавје 5 — Иден развој} идентификува пет конкретни насоки за идно истражување и развој, вклучувајќи нови карактеристики, live имплементација, пристапи базирани на Reinforcement Learning и подобрување на intel pipeline-от.

\textbf{Поглавје 6 — Заклучок} ги сумаризира главните наоди, квантитативните резултати и значењето на пристапот за полето на автоматизираното управување со спортски тимови.

\section{Технолошки стек}

Системот е целосно имплементиран во Python 3.11. Основни библиотеки и инструменти кои се користени вклучуваат: \texttt{lightgbm} за градиентно засилување, \texttt{pulp} за ILP, \texttt{anthropic} и \texttt{google-generativeai} за LLM API повици, \texttt{selenium} за веб-скрапинг, \texttt{optuna} за Bayesian оптимизација, \texttt{pandas} и \texttt{numpy} за обработка на податоци, и \texttt{flask} за корисничкиот интерфејс. Целиот изворен код е организиран во модуларна структура на директориуми со јасно разграничување меѓу фазите на пипелајнот.
